\section{结论}\label{sec:conclusion}
本文围绕体检数据连续血糖预测，保留完整评价人群，采用折内预处理和训练内部OOF选择，将基准比较、顺序优化、异质融合与固定参数预处理对照组织在同一明确协议下。原六模型中，原XGBoost在当前开发和B的MSE点估计最低，但原大网络在MAE上较低，说明模型评价需要与目标指标一致。

训练OOF选出的优化XGBoost、LightGBM、Ridge与去乙肝64--32小网络等权融合，开发/B MSE为2.177243/1.234296，相对原XGBoost下降约2.31\%/1.47\%。特征工程与单模型搜索未各自在两个集合同时改善，而有限权重融合在本轮取得较低总体平方误差；单模型较弱的网络仍可提供组合价值。

固定参数对照支持本轮两种网络采用标准化，未发现增加PCA的额外稳定收益。按训练规则改变LightGBM缺失处理后，融合开发略好、B略差，尚未显示进一步稳定提升。原大网络保留PCA是流程替换规则的结果，不代表PCA必需。

最终模型仍低估高值尾部，B集相对基线的配对差异区间跨零，且评价集合存在历史使用。因此，本文结论限于当前数据、预算和冻结配置下的复核，不将点估计改善解释为对未来样本稳定优越、临床诊断有效或效率与解释性已全面满足的证明。
